% !PRINTED 575
\weberpage{575}

%<<CONT
gen\"ugen. Diese Gleichung hat nur reelle Wurzeln, wenn $f$ gerade ist, weil dann $\tfrac{1}{2}(n - 1)$ ein Vielfaches von $e$ ist, und folglich $r$ und $r^{-1}$ in derselben Periode vorkommen. Ist aber $f$ ungerade, so sind alle Wurzeln imagin\"ar.

Es ist nun vom h\"ochsten Interesse diese Gleichung $e^{\text{ten}}$ Grades f\"ur einzelne besondere Werthe von $e$ ohne eine specielle Annahme \"uber die Primzahl $n$, ausser der, dass $n - 1$ durch $e$ theilbar sein soll, zu untersuchen.

Wir betrachten die ersten speciellen F\"alle und nehmen zun\"achst $e = 2$ an, was bei jeder ungeraden Primzahl $n$ zul\"assig ist; $\eta, \eta_1$ sind hier also die zwei Perioden von $\tfrac{1}{2}(n - 1)$ Gliedern, die wir jetzt mit $A, B$ bezeichnen wollen, so dass
\[
\begin{aligned}
A &= r + r_2 + r_4 + \cdots + r_{n-3}\\
B &= r_1 + r_3 + r_5 + \cdots + r_{n-2}.
\end{aligned}
\]
Die $r, r_2, r_4 \ldots$ haben die Exponenten $g^0, g^2 \ldots g^{n-3}$, d.\,h. die Exponenten von $g$ sind gerade Zahlen. Die Exponenten von $r$ sind also die \gesperrt{quadratischen Reste} von $n$ (\S.\,138). Ebenso sind die Exponenten von $r$ in der Summe $B$ die Nichtreste. Bezeichnen wir also die Reste mit $a$, die Nichtreste mit $b$, so ist
\weq{1}{A = \Sigma r^a, \quad B = \Sigma r^b,}
und es ist $(-1, r) = A - B$. Diese Ausdr\"ucke $A, B$, werden die \gesperrt{Gauss'schen Summen} genannt.

Machen wir in der Formel \S.\,169, (10) die Annahme
\[
e = 2, \quad \mu = f = \frac{n-1}{2}, \quad \lambda = 1, \quad \alpha = -1,
\]
so ergiebt sich
\weq{2}{A - B = \pm \sqrt{(-1)^{\frac{n-1}{2}}\, n},}
w\"ahrend andererseits
\[
A + B = -1,
\]
also
\weq{3}{2A = -1 \pm \sqrt{(-1)^{\frac{n-1}{2}}\, n}, \quad 2B = -1 \mp \sqrt{(-1)^{\frac{n-1}{2}}\, n}.}
Das Vorzeichen, das man der Wurzel zu geben hat, h\"angt von der Wahl von $r$ ab. Ist aber \"uber $r$ verf\"ugt, so ist das Vorzeichen v\"ollig bestimmt. Seine Ermittelung bietet eigenth\"umliche Schwierigkeiten, die zu einer eigenen Abhandlung von \gesperrt{Gauss} Anlass gegeben haben\textsuperscript{1}).

\noindent\rule{3cm}{0.4pt}

\noindent\textsuperscript{1}) Summatio quarundam serierum singularium. \gesperrt{Gauss} Werke, Bd.\ II. Auf andere Weise, mit Anwendung der h\"oheren Analysis, ist das Zeichen
%CONT>>
